\documentclass[11pt, a4paper]{article}

% Kodowanie znaków i język polski dla standardowego kompilatora pdfLaTeX
\usepackage[T1]{fontenc}
\usepackage[utf8]{inputenc}
\usepackage[polish]{babel}

% Czcionka bezszeryfowa
\usepackage{lmodern}
\renewcommand{\familydefault}{\sfdefault}

% Układ strony i pakiety dodatkowe
\usepackage[a4paper, margin=2.5cm]{geometry}
\usepackage{graphicx}
\usepackage{enumitem}
\setlist[itemize]{label=-}
\usepackage{amsmath}
\usepackage{amssymb}
\usepackage{booktabs}
\usepackage{tabularx}
\usepackage[normalem]{ulem}
\usepackage{hyperref}

\hypersetup{
    colorlinks=true,
    linkcolor=blue,
    urlcolor=blue
}

\setlength{\parindent}{0pt}
\setlength{\parskip}{6pt}

\title{start}
\author{Michal Pedzisz}
\date{November 2026}

\begin{document}

\maketitle

\begin{figure}[htbp]
    \centering
    \includegraphics[width=0.8\textwidth]{wykres.png}
    \caption{Podpis pod obrazkiem.}
    \label{rys:moj_wykres}
\end{figure}

\section*{Markdown, Colab oraz GitHub}

Na początek poznamy trzy narzędzia, które rozwiązują trzy różne problemy:

\begin{table}[h!]
\centering
\begin{tabularx}{\textwidth}{l X}
\toprule
\textbf{Narzędzie} & \textbf{Do czego służy?} \\
\midrule
\textbf{Markdown} & do prostego zapisywania uporządkowanego tekstu, dokumentacji i raportów \\
\textbf{Google Colab} & do łączenia tekstu, kodu, obliczeń i wyników w jednym notebooku \\
\textbf{GitHub} & do przechowywania projektu, śledzenia zmian i współpracy \\
\bottomrule
\end{tabularx}
\end{table}

Te trzy elementy bardzo dobrze ze sobą współpracują. Podstawą wszędzie jest przygotowanie opisu w kodzie Markdown. Do tego dochodzić będą obliczenia w Colabie, a następnie przechowanie źródła i wyniki w repozytorium GitHub.

\medskip\hrule\medskip

\section{Markdown}

\subsection{Co to jest Markdown?}

\textbf{Markdown} to prosty język znaczników służący do formatowania zwykłego tekstu.

Najczęściej zapisujemy go w plikach z rozszerzeniem:

\begin{verbatim}
.md
\end{verbatim}

Przykłady:

\begin{verbatim}
README.md
raport.md
notatki.md
\end{verbatim}

Plik Markdown jest zwykłym plikiem tekstowym. Można go otworzyć praktycznie w każdym edytorze. Specjalne znaki, takie jak \texttt{\#}, \texttt{*}, \texttt{-} czy \texttt{\$}, informują program wyświetlający dokument, jak ma wyglądać tekst.

\section{Dlaczego Markdown jest tak istotny?}

Za każdym razem rozmawiając z chatem AI tak naprawdę rozmawiamy z plikiem Markdown. AI potrafi interpretować składnię Markdown i generować raporty w tym formacie.

Dlatego też jak skopiujesz output to ma on całą masę znaczników (np. \texttt{\#}, \texttt{*}, \texttt{-}, czy \texttt{\$}), które są potrzebne do poprawnego renderowania.

Github (archiwum programistyczne) zawiera opisy również w Markdown. Podobnie Google Colab (środowisko programistyczne pythona) potrafią renderować Markdown. Dzięki temu możemy przygotować raport, który będzie wyglądał dobrze zarówno w przeglądarce, jak i w notebooku.

Warto od razu rozróżnić dwie rzeczy:

\begin{itemize}
    \item \textbf{źródło Markdown} --- tekst, który wpisujemy do pliku,
    \item \textbf{wyrenderowany dokument} --- wygląd tego tekstu po interpretacji składni Markdown.
\end{itemize}

Na GitHubie możemy przełączać się między widokiem źródła i wyrenderowanym dokumentem.

Podobnie edytory kodu jak VS Code, PyCharm czy Jupyter Notebook potrafią renderować Markdown w osobnym oknie podglądu.

\subsection{Zanim zaczniemy pisać w Markdown zobacz przykładowy plik Markdown}

Raport na temat ciała na sprężynie:

\begin{itemize}
    \item \href{examples/01_markdown/raport_cialo_na_sprezynie.md}{,,Ciało na sprężynie''}.
\end{itemize}

\subsection{Nagłówki}

Nagłówki budują strukturę dokumentu. Im więcej znaków \texttt{\#}, tym niższy poziom nagłówka.

Źródło:

\begin{verbatim}
# Tytuł dokumentu

## Główna sekcja

### Podsekcja
\end{verbatim}

Efekt:

\begin{quote}
{\LARGE \textbf{Tytuł dokumentu}}\par\medskip
{\Large \textbf{Główna sekcja}}\par\medskip
{\large \textbf{Podsekcja}}
\end{quote}

W praktyce jeden dokument powinien mieć jeden główny nagłówek \texttt{\#}.

Dalsze części dokumentu organizujemy za pomocą \texttt{\#\#}, \texttt{\#\#\#} itd.

\medskip\hrule\medskip

\subsection{Pogrubienie, kursywa i przekreślenie}

Źródło:

\begin{verbatim}
**tekst pogrubiony**

*tekst pochylony*

~~tekst przekreślony~~

***tekst pogrubiony i pochylony***
\end{verbatim}

Efekt:

\textbf{tekst pogrubiony}

\textit{tekst pochylony}

\sout{tekst przekreślony}

\textbf{\textit{tekst pogrubiony i pochylony}}

Wyróżnień warto używać oszczędnie. Jeśli wszystko jest pogrubione, nic nie jest naprawdę wyróżnione.
\medskip\hrule\medskip

\subsection{Akapit i nowa linia}

Nowy akapit tworzymy przez pozostawienie pustej linii.

Źródło:

\begin{verbatim}
To jest pierwszy akapit.

To jest drugi akapit.
\end{verbatim}

Efekt:

To jest pierwszy akapit.

To jest drugi akapit.

To ważna zasada: w Markdown pojedyncze naciśnięcie Enter nie zawsze oznacza nowy akapit.

\medskip\hrule\medskip

\subsection{Listy punktowane}

Źródło:

\begin{verbatim}
- Markdown
- Colab
- GitHub
\end{verbatim}

Efekt:

\begin{itemize}
    \item Markdown
    \item Colab
    \item GitHub
\end{itemize}

Listy można zagnieżdżać:

\begin{verbatim}
- GitHub
  - repozytorium
  - Issues
  - historia zmian
- Colab
  - tekst
  - kod
  - wyniki
\end{verbatim}

Efekt:

\begin{itemize}
    \item GitHub
    \begin{itemize}
        \item repozytorium
        \item Issues
        \item historia zmian
    \end{itemize}
    \item Colab
    \begin{itemize}
        \item tekst
        \item kod
        \item wyniki
    \end{itemize}
\end{itemize}

\medskip\hrule\medskip

\subsection{Listy numerowane}

Źródło:

\begin{verbatim}
1. przygotuj dane,
2. wykonaj obliczenia,
3. zapisz wynik,
4. opisz wynik.
\end{verbatim}

Efekt:

\begin{enumerate}
    \item przygotuj dane,
    \item wykonaj obliczenia,
    \item zapisz wynik,
    \item opisz wynik.
\end{enumerate}

Listy numerowane są szczególnie przydatne przy opisie procedury lub instrukcji.

\medskip\hrule\medskip

\subsection{Checklisty}

GitHub obsługuje listy zadań.

Źródło:

\begin{verbatim}
- [x] utworzono notebook
- [x] wykonano obliczenia
- [ ] sprawdzono wynik
\end{verbatim}

Efekt:

\begin{itemize}
    \item[$\boxtimes$] utworzono notebook
    \item[$\boxtimes$] wykonano obliczenia
    \item[$\square$] sprawdzono wynik
\end{itemize}

Checklisty bardzo dobrze nadają się do krótkich list kontrolnych w dokumentacji i GitHub Issues.

\medskip\hrule\medskip

\subsection{Linki}

Źródło:

\begin{verbatim}
[GitHub](https://github.com/)
\end{verbatim}

Efekt:

\href{https://github.com/}{GitHub}

Zamiast wklejać długi adres do strony, warto umieścić go pod czytelną nazwą.

\medskip\hrule\medskip

\subsection{Obrazy}

Obraz osadzamy podobnie jak link, ale na początku dodajemy znak \texttt{!}.

Źródło:

\begin{verbatim}
![Wykres zależności y od x](wykres.png)
\end{verbatim}

Jeżeli plik \texttt{wykres.png} znajduje się w tym samym folderze co dokument Markdown, GitHub wyświetli obraz w miejscu tej instrukcji.

Jeżeli plik znajduje się w podfolderze \texttt{files}, używamy ścieżki:

\begin{verbatim}
![Wykres zależności y od x](files/wykres.png)
\end{verbatim}

Efekt:

\framebox{\parbox{0.8\textwidth}{\centering Wykres zależności y od x (\texttt{files/wykres.png})}}

Tekst w nawiasach kwadratowych jest opisem obrazu. Warto go uzupełniać --- pomaga zrozumieć zawartość także wtedy, gdy obraz nie może zostać wyświetlony.

W projektach przechowywanych w repozytorium zwykle wygodniej używać \textbf{ścieżek względnych} niż pełnych adresów internetowych.

\medskip\hrule\medskip

\subsection{Tabele}

Źródło:

\begin{verbatim}
| Narzędzie | Typ pliku | Zastosowanie |
| --- | --- | --- |
| Markdown | .md | dokumentacja |
| Colab | .ipynb | obliczenia |
| GitHub | repozytorium | historia projektu |
\end{verbatim}

Efekt:

\begin{table}[h!]
\centering
\begin{tabular}{lll}
\toprule
\textbf{Narzędzie} & \textbf{Typ pliku} & \textbf{Zastosowanie} \\
\midrule
Markdown & \texttt{.md} & dokumentacja \\
Colab & \texttt{.ipynb} & obliczenia \\
GitHub & repozytorium & historia projektu \\
\bottomrule
\end{tabular}
\end{table}

Tabele są dobre do krótkich zestawień. Jeśli w komórkach pojawiają się długie akapity tekstu, zwykle lepiej wrócić do zwykłych sekcji i list.

\medskip\hrule\medskip

\subsection{Kod w tekście}

Nazwy plików, komendy i bardzo krótkie fragmenty kodu warto oznaczać pojedynczym backtickiem.

Źródło:

\begin{verbatim}
Uruchom plik `analiza.py` i sprawdź wartość zmiennej `wynik`.
\end{verbatim}

Efekt:

Uruchom plik \texttt{analiza.py} i sprawdź wartość zmiennej \texttt{wynik}.

\medskip\hrule\medskip

\subsection{Bloki kodu}

Większy fragment kodu umieszczamy w osobnym bloku.

Źródło:

\begin{verbatim}
```python
x = 5
y = x**2
print(y)
\end{verbatim}

Podanie nazwy języka, np. \texttt{python}, \texttt{bash}, \texttt{html} czy \texttt{json}, pozwala GitHubowi kolorować składnię.

W raporcie technicznym dobrze jest pokazać krótki fragment kodu, który jest istotny dla opisywanego wyniku. Nie trzeba kopiować do raportu całego programu.
```latex
\medskip\hrule\medskip

\subsection{Cytaty}

Źródło:

\begin{verbatim}
> Wynik należy zawsze sprawdzić przed jego wykorzystaniem.
\end{verbatim}

Efekt:

\begin{quote}
Wynik należy zawsze sprawdzić przed jego wykorzystaniem.
\end{quote}

Cytaty można wykorzystać także do krótkiego wyróżnienia komentarza lub uwagi.

\medskip\hrule\medskip

\subsection{Matematyka}

GitHub potrafi renderować wzory matematyczne.

Krótki wzór umieszczamy wewnątrz pojedynczych znaków dolara.

Źródło:

\begin{verbatim}
Energia spoczynkowa dana jest wzorem $E=mc^2$.
\end{verbatim}

Efekt:

Energia spoczynkowa dana jest wzorem $E=mc^2$.

Dłuższy wzór umieszczamy w osobnym bloku.

Źródło:

\begin{verbatim}
$$
\sum_{k=1}^{n} k^2 = \frac{n(n+1)(2n+1)}{6}
$$
\end{verbatim}

Efekt:

\[
\sum_{k=1}^{n} k^2 = \frac{n(n+1)(2n+1)}{6}
\]

Macierze zapisujemy wielowierszowo:

\begin{verbatim}
$$
A=
\begin{pmatrix}
1 & 2 \\
3 & 4 \\
\end{pmatrix}
$$
\end{verbatim}

Efekt:

\[
A=
\begin{pmatrix}
1 & 2 \\
3 & 4 \\
\end{pmatrix}
\]

\medskip\hrule\medskip

\subsection{Zrzut ekranu i transkrypcja do Markdown}

Czasem informacja istnieje tylko jako obraz: zrzut ekranu programu, zdjęcie tabeli, komunikat błędu albo fragment zeskanowanego dokumentu. AI może pomóc przepisać zawartość obrazu do tekstu. Taki proces nazywamy tutaj \textbf{transkrypcją}.

Przykładowy workflow:

\begin{enumerate}
    \item zachowujemy oryginalny obraz,
    \item prosimy AI o \textbf{przepisanie jego zawartości do kodu markdown},
    \item zapisujemy wynik w pliku tekstowym lub Markdown,
    \item porównujemy transkrypcję z obrazem,
    \item poprawiamy ewentualne błędy.
\end{enumerate}

Największej ostrożności wymagają:

\begin{itemize}
    \item liczby,
    \item wzory,
    \item nazwy plików,
    \item adresy,
    \item fragmenty kodu,
    \item komunikaty błędów.
\end{itemize}

Transkrypcja wygenerowana przez AI jest \textbf{wersją roboczą}, dopóki człowiek jej nie sprawdzi.

\medskip\hrule\medskip

\subsection{Poprawne formatowanie matematyki w Markdown}

W materiałach tego kursu polecamy stosowanie prostego zestawu reguł formatowania matematyki przygotowanego specjalnie tak, aby pliki dobrze renderowały się na GitHubie:

\begin{itemize}
    \item \href{hints.md}{hints.md}
\end{itemize}

Pracując z chatem AI i domagając się kodu markdowna, warto załączać ten plik jako kontekst, aby AI od razu generowało poprawny kod w Markdown.

Różne systemy: przeglądarkowe chaty AI, Edytor Visual Studio, Colab i GitHub, mogą różnie renderować wzory matematyczne!

Częstym błędem jest używanie w kodzie Markdown znaków \texttt{\textbackslash[]} i \texttt{\textbackslash]} zamiast \texttt{\$ \$} do oznaczania wzorów matematycznych. Wtedy GitHub nie renderuje wzoru, a Colab renderuje go poprawnie.

\medskip\hrule\medskip

\subsection{Dobre materiały o Markdown}

\begin{itemize}
    \item \href{https://docs.github.com/en/get-started/writing-on-github/getting-started-with-writing-and-formatting-on-github/basic-writing-and-formatting-syntax}{GitHub Docs --- Basic writing and formatting syntax}
    \item \href{https://docs.github.com/en/get-started/writing-on-github/working-with-advanced-formatting/writing-mathematical-expressions}{GitHub Docs --- Writing mathematical expressions}
    \item \href{https://www.markdownguide.org/basic-syntax/}{Markdown Guide --- Basic Syntax}
    \item \href{https://www.markdownguide.org/cheat-sheet/}{Markdown Guide --- Cheat Sheet}
\end{itemize}

\end{document}